\theory{Group theory}{What can be learned from studying reversible operations and the symmetries they form?}{A group is a collection of reversible actions that can be composed, with an identity action and an inverse for every action. Group theory studies the common structure behind rotations, permutations, arithmetic symmetries, and transformations.}{Geometry, algebra, number theory, physics, chemistry, cryptography, coding, and the classification of finite structures.}{Composition and inverses encode symmetry operations, while subgroups and representations expose the internal structure of those symmetries.}

\paragraph{The idea of a group}
Imagine rotating a square around its center. There are four rotations that return the square to itself: turn by $0°$, $90°$, $180°$, or $270°$. Doing one rotation after another gives another rotation (for example, $90°$ then $90°$ gives $180°$). Doing nothing is an action too (the identity). Every rotation can be undone: to undo a $90°$ turn, rotate by $270°$. These four rules---combining actions, an identity, and an inverse for each action---define a group. The elements need not be numbers; they may be rotations, shuffles, reflections, or any reversible operation.

Formally, a group has a binary operation (a way to combine two actions into one), an identity element (doing nothing), and an inverse for every element (a way to undo). A group is abelian if the order of composition does not matter: rotating $90°$ then $180°$ gives the same result as $180°$ then $90°$. A group that is not abelian---like the symmetries of a square including reflections---is called non-abelian \cite{group_theory_ref1}.

\end{historylist}
\begin{quote}\small \textbf{Prerequisite:} High-school algebra. The first two sections need no proof techniques.\end{quote}
\section*{3. Theory}
\subsection*{Core theory}
\paragraph{Subgroups, cosets, and quotients}
A subgroup is a smaller collection closed under the same operation. The even integers form a subgroup of the integers under addition. A subgroup partitions a group into equal-sized left or right cosets. Lagrange's theorem follows: the order of a finite subgroup divides the order of the whole group.

If a subgroup is normal, its cosets can themselves be multiplied to form a quotient group. Normal subgroups are precisely the subgroups compatible with conjugation by every group element. A homomorphism is a structure-preserving map between groups; its kernel is normal, and the first isomorphism theorem says that the quotient by the kernel is isomorphic to the image.

\paragraph{How groups solve symmetry problems}
The symmetries of an object act on the object, producing a group action. The orbit of a point is the collection of positions reachable by the group, while the stabilizer consists of actions that leave it fixed. Orbit-stabilizer says, for a finite group,
\[
|G|=|\operatorname{Orb}(x)|\,|\operatorname{Stab}(x)|.
\]
This counts symmetries by counting where an object can go and how many symmetries keep it there.

Permutation groups rearrange a finite set. The symmetric group $S_n$ contains all $n!$ permutations; the alternating group $A_n$ contains the even permutations. Cayley's theorem embeds every group into a permutation group, showing that permutations are a universal model for abstract groups.

\paragraph{Finite groups and their classification}
Finite groups may be described by a multiplication table, generators and relations, or a presentation. Simple groups have no nontrivial normal subgroup. They are the indivisible building blocks of finite groups, in the same way that prime numbers are building blocks of integers. The classification of finite simple groups is a vast theorem showing that every finite simple group belongs to one of several families: cyclic groups of prime order, alternating groups, groups of Lie type, or one of twenty-six sporadic groups. The proof occupies thousands of pages \cite{group_theory_ref2}.

The Sylow theorems constrain subgroups whose order is a power of a prime. They often prove that a group of a proposed order must be abelian or must contain a normal subgroup. Semidirect products explain how one group can act on another to build a larger group.

\paragraph{Representations and characters}
A representation realizes abstract group elements as matrices acting on a vector space. Matrix multiplication then performs the group operation, so linear algebra becomes available. A representation is faithful when different group elements give different matrices. Characters record traces of representation matrices; they are constant on conjugacy classes and can determine how representations decompose.

Fourier polynomials are characters of the circle group $U(1)$ acting on periodic functions. In quantum mechanics, a state space carries a representation of a symmetry group, and conserved quantities arise from symmetry through the broader principle associated with Noether. Finite-group representations are also used to analyze molecules and crystal symmetries \cite{group_theory_ref3}.

\paragraph{Lie groups and continuous symmetry}
A Lie group is both a group and a differentiable manifold, with smooth multiplication and inversion. Rotations in three-dimensional space form the Lie group $SO(3)$. Its tangent space at the identity is a Lie algebra, whose bracket measures the first-order failure of rotations to commute.

Lie groups are the natural framework for continuous symmetries of differential equations, geometry, and physics. Differential Galois theory studies continuous symmetries of differential equations in a role analogous to permutation groups for algebraic equations. Compact Lie groups and their representations describe rotations, angular momentum, and gauge symmetries.

\paragraph{Combinatorial and geometric group theory}
Combinatorial group theory studies presentations, words in generators, and the word problem: can an algorithm decide whether a given word represents the identity? Free groups have no relations beyond cancellation. Adding relations produces quotient groups and often a geometric space.

Geometric group theory studies a group through a space on which it acts. A Cayley graph records generators as labeled edges, turning algebraic multiplication into paths. Hyperbolic groups imitate negative-curvature geometry, while Bass--Serre theory studies groups acting on trees. These ideas connect group theory to topology and geometry.

\paragraph{Applications}
Elliptic curves over finite fields have groups of points. Their discrete logarithm problem supports public-key cryptography. Diffie--Hellman key exchange uses a finite cyclic group; braid-group and other non-abelian proposals use harder group operations, though security must be analyzed carefully. Error-correcting codes use group structure to organize words and decode them. Caesar's cipher is a simple group operation on alphabet positions \cite{group_theory_ref4}.

Symmetry groups classify crystal patterns, molecules, tilings, and physical laws. In topology, fundamental groups record loops and coverings. In number theory, Galois groups encode symmetries of field extensions. In algebraic geometry, algebraic groups and group schemes add a compatible group operation to a geometric object.

\paragraph{What the theory depends on and where it is useful}
Group theory depends on set operations, functions, relations, and increasingly on linear algebra, topology, and geometry. It helps Galois theory organize root symmetries, representation theory convert algebra to matrices, Lie theory describe continuous motion, topology classify loops, and cryptography construct one-way computational tasks.

The classification of a group may be difficult even when its definition is short. A group model also records exact symmetry, not measurement noise or approximate symmetry; applications often require additional geometry, probability, or computation.

\section*{4. Important theorems and results}
\begin{enumerate}[leftmargin=*,itemsep=0.35em]
\item \textbf{Lagrange's theorem.} The order of a subgroup of a finite group divides the order of the group. \emph{Example.} $S_3$ has order $6$; its subgroups have orders $1,2,3$, and $6$, all divisors of $6$.

\item \textbf{Cayley's theorem.} Every group is isomorphic to a subgroup of a symmetric permutation group. \emph{Example.} The cyclic group $\mathbb{Z}/3\mathbb{Z}$ embeds in $S_3$ by sending the generator $1$ to the 3-cycle $(1\;2\;3)$.

\item \textbf{Orbit-stabilizer theorem.} For a finite group action, $|G|=|\operatorname{Orb}(x)|\,|\operatorname{Stab}(x)|$. \emph{Example.} $S_3$ acts on $\{1,2,3\}$. The orbit of $1$ is $\{1,2,3\}$ (size 3). The stabilizer of $1$ is $\{\text{id},(2\;3)\}$ (size 2). Then $|S_3|=3\cdot 2=6$.

\item \textbf{First isomorphism theorem.} A quotient by the kernel of a homomorphism is isomorphic to its image. \emph{Example.} The map $\varphi:\mathbb{Z}\to\mathbb{Z}/6\mathbb{Z}$ has kernel $6\mathbb{Z}$, so $\mathbb{Z}/6\mathbb{Z}\cong\mathbb{Z}/6\mathbb{Z}$, confirming the theorem.

\item \textbf{Classification of finite simple groups.} Every finite simple group belongs to the cyclic, alternating, Lie-type, or sporadic families \cite{group_theory_ref2}.

These results move from local structure to classification. Lagrange and orbit--stabilizer give numerical tests for a finite action, Cayley realizes every group as permutations, and the isomorphism theorem explains how kernels remove information. The classification theorem is a map of the possible indivisible finite symmetries; it is not a short method for identifying an arbitrary group.
\end{enumerate}

\theoryapplications
\theoryconnections{group theory}{%
\item Set theory supplies the underlying collection and operation, while algebra studies homomorphisms, subgroups, quotients, and actions.
\item Geometry and number theory provide important examples, and representation theory turns abstract group elements into matrices.
}{%
\item Group theory helps describe symmetry in geometry, chemistry, physics, coding, and algorithms.
\item It also organizes transformations into reusable algebraic objects, so a proof about one group action can apply to many different systems.
}
\section*{Glossary}
\begin{description}[leftmargin=*,style=nextline,itemsep=0.3em]
\item[\textbf{Group.}] A set with an associative binary operation, an identity element, and an inverse for every element; it formalizes the idea of symmetry or reversible transformations.
\item[\textbf{Subgroup.}] A subset of a group that is itself a group under the same operation.
\item[\textbf{Normal subgroup.}] A subgroup invariant under conjugation by every group element, which allows the formation of a quotient group.
\item[\textbf{Homomorphism.}] A structure-preserving map between two groups; its kernel is always a normal subgroup of the domain.
\item[\textbf{Orbit.}] The set of positions reachable from a given element by applying all elements of an acting group.
\item[\textbf{Simple group.}] A group with no nontrivial normal subgroup; simple groups are the building blocks from which all finite groups are assembled.
\item[\textbf{Representation.}] A realization of an abstract group as matrices acting on a vector space, making linear-algebraic tools available for study.
\item[\textbf{Coset.}] A set $gH=\{gh:h\in H\}$ for a subgroup $H$ and element $g$; cosets partition the group into equal-sized pieces.
\item[\textbf{Quotient group.}] The group $G/N$ formed by cosets of a normal subgroup $N$.
\item[\textbf{Kernel.}] The set of elements mapped to the identity by a homomorphism; always a normal subgroup.
\item[\textbf{Group action.}] A homomorphism from a group to the permutation group of a set.
\item[\textbf{Abelian group.}] A group in which the operation is commutative.
\end{description}
\section*{7. Sample Questions}
\emph{Exercise reference:} \cite{group_theory_ref1}. The cited edition is the source for the exercise set; consult its Exercises section for the official statement and numbering.
\begin{enumerate}[leftmargin=*,itemsep=0.9em]
\item \textbf{Exercise 1.} Find all subgroups of $\mathbb{Z}/12\mathbb{Z}$ and verify Lagrange's theorem for each.

\emph{Solution.} By the correspondence theorem, subgroups of $\mathbb{Z}/12\mathbb{Z}$ are $\langle d\rangle$ for $d\mid 12$: $\langle 0\rangle=\{0\}$ (order 1), $\langle 1\rangle=\mathbb{Z}/12\mathbb{Z}$ (order 12), $\langle 2\rangle=\{0,2,4,6,8,10\}$ (order 6), $\langle 3\rangle=\{0,3,6,9\}$ (order 4), $\langle 4\rangle=\{0,4,8\}$ (order 3), $\langle 6\rangle=\{0,6\}$ (order 2). Each subgroup order divides 12, confirming Lagrange's theorem.

\item \textbf{Exercise 2.} Compute the order of the element $(1\;2\;3)$ in $S_4$.

\emph{Solution.} The permutation $(1\;2\;3)$ is a 3-cycle. Its order is $3$, since $(1\;2\;3)^3=\text{id}$ and no smaller positive power is the identity: $(1\;2\;3)^1\neq\text{id}$, $(1\;2\;3)^2=(1\;3\;2)\neq\text{id}$.

\item \textbf{Exercise 3.} Find the kernel of the homomorphism $\varphi:\mathbb{Z}\to\mathbb{Z}/6\mathbb{Z}$ given by $\varphi(n)=n\bmod 6$.

\emph{Solution.} $\ker\varphi=\{n\in\mathbb{Z}:n\equiv 0\pmod{6}\}=6\mathbb{Z}$. By the first isomorphism theorem, $\mathbb{Z}/\ker\varphi=\mathbb{Z}/6\mathbb{Z}\cong\operatorname{im}\varphi=\mathbb{Z}/6\mathbb{Z}$, which checks out.

\item \textbf{Exercise 4.} Prove that a group of order $7$ is cyclic.

\emph{Solution.} By Lagrange's theorem, every non-identity element has order dividing $7$. Since $7$ is prime, the only divisors are $1$ and $7$. Any non-identity element has order $7$, hence generates the entire group. Therefore the group is cyclic: $G\cong\mathbb{Z}/7\mathbb{Z}$.

\item \textbf{Exercise 5.} Determine whether the subgroup $\{(1),(1\;2)(3\;4),(1\;3)(2\;4),(1\;4)(2\;3)\}$ of $S_4$ is normal.

\emph{Solution.} This subgroup $V\cong\mathbb{Z}/2\times\mathbb{Z}/2$ (the Klein four-group) is normal in $S_4$. To see why, observe that $V$ consists of the identity and all products of two disjoint transpositions. Conjugation by any $\sigma\in S_4$ sends a product of two disjoint transpositions to another product of two disjoint transpositions (it simply relabels the elements), so $\sigma V\sigma^{-1}\subseteq V$. Since $|V|$ is fixed and conjugation is a bijection, $\sigma V\sigma^{-1}=V$ for every $\sigma\in S_4$.

\end{enumerate}
\section*{8. References}
\addcontentsline{toc}{section}{References}

\begin{thebibliography}{9}
\bibitem{group_theory_ref1} Joseph J. Rotman, \emph{An Introduction to the Theory of Groups}, 4th ed., Springer, 1995.
\bibitem{group_theory_ref2} Daniel Gorenstein, Richard Lyons, and Ronald Solomon, \emph{The Classification of the Finite Simple Groups}, American Mathematical Society, 1994--.
\bibitem{group_theory_ref3} Jean-Pierre Serre, \emph{Linear Representations of Finite Groups}, Springer, 1977.
\bibitem{group_theory_ref4} David S. Dummit and Richard M. Foote, \emph{Abstract Algebra}, 3rd ed., Wiley, 2004.
\end{thebibliography}
